\documentclass[11pt]{amsart}

\usepackage[T1]{fontenc}
\usepackage{lmodern}
\usepackage{microtype}
\usepackage{amsmath,amssymb,amsthm}
\usepackage{booktabs}
\usepackage[hidelinks]{hyperref}


\hypersetup{
  pdftitle={The Least Counterexample to Conjecture 4.3 of Dolfi, Hafezieh, and Spiga}
}

\newtheorem{theorem}{Theorem}
\newtheorem{lemma}[theorem]{Lemma}
\newtheorem{proposition}[theorem]{Proposition}
\theoremstyle{remark}
\newtheorem{remark}[theorem]{Remark}

\DeclareMathOperator{\ord}{ord}

\title[Least Counterexample to Conjecture 4.3]
{The Least Counterexample to Conjecture 4.3 of Dolfi, Hafezieh, and Spiga}

\date{}

\subjclass[2020]{11A51; 11A41, 20C15}
\keywords{Mersenne number, cyclotomic polynomial, prime divisor,
counterexample}

\begin{document}

\begin{abstract}
Dolfi, Hafezieh, and Spiga conjectured that if a positive integer $n$
has $\ell\ge 3$ distinct prime divisors and $\gcd(2^n-1,n)=1$, then
$2^n-1$ has at least $2^\ell$ distinct prime divisors. We observe that
the conjecture is false as stated and prove that its least counterexample
is $n=170$:
\[
\gcd(2^{170}-1,170)=1,
\qquad
\omega(2^{170}-1)=7<8=2^{\omega(170)}.
\]
The factorization itself appears in prior data; the point of the note is
to identify its consequence for the conjecture and prove minimality.
Since $170$ is even, this does not address the odd-exponent case
motivating the original group-theoretic construction.
\end{abstract}

\maketitle

\section{The counterexample}

For a positive integer $m$, let $\omega(m)$ be the number of distinct
prime divisors of $m$, and let $\tau(m)$ be the number of positive
divisors of $m$. Conjecture~4.3 of Dolfi, Hafezieh, and Spiga
\cite{DHS} reads, in full:
\begin{quote}
``Let $n$ be divisible by $\ell$ distinct primes and assume
$\gcd(2^n-1,n)=1$. If $\ell\ge3$, then $2^n-1$ is divisible by at
least $2^\ell$ distinct primes.''
\end{quote}
That is,
\[
\gcd(2^n-1,n)=1,
\quad \omega(n)=\ell\ge 3
\quad\Longrightarrow\quad
\omega(2^n-1)\ge 2^\ell,
\]
with no hypothesis of oddness on $n$. Immediately before the
conjecture the same authors write: ``Except when $\ell\le2$, we have no
examples where $2^n-1$ has exactly $2^{|\pi(n)|}-1$ distinct prime
divisors.'' Here $|\pi(n)|=\omega(n)$, and the exponent exhibited below
is such an example.

\begin{theorem}\label{thm:main}
The least positive integer violating Conjecture~4.3 of \cite{DHS} is
\[
n=170=2\cdot5\cdot17.
\]
More precisely,
\[
\gcd(2^{170}-1,170)=1
\quad\text{and}\quad
\omega(2^{170}-1)=7.
\]
\end{theorem}

We use the following specialization of \cite[Lemma~4.2]{DHS}.

\begin{lemma}\label{lem:cyclotomic}
If $\gcd(2^n-1,n)=1$, then
\[
\omega(2^n-1)\ge \tau(n)-1.
\]
\end{lemma}

\begin{proof}
For every divisor $d>1$ of $n$, one has $\Phi_d(2)>1$; choose a prime
$r_d\mid\Phi_d(2)$. Then $r_d\mid2^n-1$, so the gcd hypothesis gives
$r_d\nmid n$, and hence $r_d\nmid d$. We claim that
$\ord_{r_d}(2)=d$. Since $r_d\mid\Phi_d(2)\mid 2^d-1$, the order
$o=\ord_{r_d}(2)$ divides $d$; suppose $o<d$. Modulo $r_d$, the element
$2$ is a root of $x^o-1=\prod_{e\mid o}\Phi_e(x)$ and therefore of some
$\Phi_e(x)$ with $e\mid o$, and $e\le o<d$ forces $e\ne d$. Thus $e$
and $d$ are distinct divisors of $d$, so $\Phi_e\Phi_d$ divides
$x^d-1$, and $2$ is a repeated root of $x^d-1$ over
$\mathbf F_{r_d}$. This is impossible because $x^d-1$ is squarefree
over $\mathbf F_{r_d}$ when $r_d\nmid d$. Therefore the primes $r_d$
are distinct as $d$ varies, giving $\tau(n)-1$ distinct prime divisors
of $2^n-1$.
\end{proof}

\begin{proof}[Proof of Theorem~\ref{thm:main}]
Direct evaluation of the cyclotomic factors for the divisors of $170$
gives
\begin{equation}\label{eq:factorization}
\begin{split}
2^{170}-1={}&3\cdot11\cdot31\cdot43691\cdot131071\\
&{}\cdot9520972806333758431
\cdot26831423036065352611.
\end{split}
\end{equation}
The first five factors are prime by trial division, and
Proposition~\ref{prop:large-primes} gives deterministic primality
certificates for the last two. Thus $\omega(2^{170}-1)=7$.

Moreover, $2^{170}-1$ is odd. Since $\ord_5(2)=4$ and
$\ord_{17}(2)=8$, while neither $4$ nor $8$ divides $170$, neither $5$
nor $17$ divides $2^{170}-1$. Hence
\[
\gcd(2^{170}-1,170)=1,
\]
so $170$ is a counterexample.

It remains to prove minimality. Suppose that $n<170$ satisfies the gcd
hypothesis, has $\ell=\omega(n)\ge3$, and violates the conjectured bound.
If $n$ is not squarefree, then
\[
\tau(n)\ge3\cdot2^{\ell-1},
\]
and Lemma~\ref{lem:cyclotomic} yields
\[
\omega(2^n-1)\ge\tau(n)-1
\ge3\cdot2^{\ell-1}-1\ge2^\ell,
\]
a contradiction. Thus $n$ is squarefree. Since
$2\cdot3\cdot5\cdot7=210$, we must have $\ell=3$.

The squarefree products of three distinct primes below $170$ are
\[
30,42,66,70,78,102,105,110,114,130,138,154,165.
\]
Every value divisible by $6$ fails the gcd hypothesis because
$3\mid n$ and, $n$ being even, $3\mid 2^n-1$ as well; hence
$3\mid\gcd(2^n-1,n)$. The value $105$ fails because
$\ord_7(2)=3\mid105$, and $110$ fails because
$\ord_{11}(2)=10\mid110$. Hence any counterexample below $170$ must lie
in
\[
\{70,130,154,165\}.
\]
For each remaining value, the proof of Lemma~\ref{lem:cyclotomic}
shows that every prime divisor of $\Phi_d(2)$ has order $d$; each such
$n$ is squarefree with three prime factors, so $\tau(n)=8$, and hence
the seven nontrivial cyclotomic factors contribute disjoint sets of
primes. The following factorizations show that one such set has at
least two members:
\[
\begin{array}{c|c}
 n & \text{a factored cyclotomic value }\Phi_d(2),\ d\mid n\\
\hline
70  & \Phi_{35}(2)=71\cdot122921,\\
130 & \Phi_{130}(2)=131\cdot409891\cdot7623851,\\
154,165 & \Phi_{11}(2)=23\cdot89.
\end{array}
\]
All displayed factors are prime by trial division. Thus
$\omega(2^n-1)\ge8$ in every remaining case, completing the proof.
\end{proof}

\begin{remark}[Scope, source of the wording, and prior data]
\label{rem:scope}
The counterexample concerns the unrestricted wording of
Conjecture~4.3 quoted in Section~1. That wording is the one printed in
the preprint version of \cite{DHS}, which is the version we were able
to consult; the statement number is corroborated for the version of
record only indirectly, by \cite{LLRS}, who cite ``Conjecture 4.3 on
page 351''. If the published version carries an additional hypothesis,
then what Theorem~\ref{thm:main} refutes is the sentence displayed in
Section~1.

The construction in Section~3 of \cite{DHS} assumes more than the
conjecture does: it opens ``let $n$ be an odd positive integer such
that $n\ne p$, $n_3=3$ and $\gcd(p^n-1,n)=1$''. Since $170$ is even
and $3\nmid170$, Theorem~\ref{thm:main} says nothing about that
regime, and the group-theoretic construction resting on it is
untouched.

The factorization itself is not new. It is classical
Cunningham-project data: $2^{170}-1=(2^{85}-1)(2^{85}+1)$, and both
factors are recorded in the tables of Brillhart, Lehmer, Selfridge,
Tuckerman, and Wagstaff \cite{BLSTW}, which OEIS A046800 also lists
among its references; A046800 records $\omega(2^{170}-1)=7$
\cite{OEIS}.
Example~3.8 of Laubacher, Lewis, Ravaglia, and Summers \cite{LLRS}
records the six factors in \eqref{eq:factorization} other than $3$ by
taking $p_1=5$ and $p_2=17$, and cites Conjecture~4.3 on the following
page without connecting the two, their construction being confined to
odd exponents. No novelty is claimed for the factorization; the point
of the note is its consequence for Conjecture~4.3 and the proof of
minimality.
\end{remark}

\section{Primality certificates}

We include short Pocklington certificates for the two large factors in
\eqref{eq:factorization}.

\begin{lemma}[Pocklington criterion]\label{lem:pocklington}
Let $N>1$ be odd, and let $F$ be a squarefree divisor of $N-1$ whose
prime factorization is known. Suppose that $F^2>N$ and that an integer
$a$ satisfies
\[
a^{N-1}\equiv1\pmod N,
\qquad
\gcd\bigl(a^{(N-1)/q}-1,N\bigr)=1
\]
for every prime $q\mid F$. Then $N$ is prime.
\end{lemma}

\begin{proof}
Let $r$ be a prime divisor of $N$. For every $q\mid F$, the displayed
conditions imply that $q\mid\ord_r(a)$, and hence $q\mid r-1$.
Therefore $F\mid r-1$, so $r>F>\sqrt N$. A composite $N$ has a prime
divisor at most $\sqrt N$, a contradiction.
\end{proof}

\begin{proposition}\label{prop:large-primes}
The integers
\[
P=9520972806333758431,
\qquad
Q=26831423036065352611
\]
are prime.
\end{proposition}

\begin{proof}
The integers
\[
709,\ 2437,\ 2857,\ 6367,\ 27953
\]
are prime by trial division. In each row of the following table,
$F\mid N-1$, $F^2>N$, and the displayed witness $a$ satisfies the two
conditions in Lemma~\ref{lem:pocklington} for every prime $q\mid F$.
The rows are read from top to bottom within each of the two certificate
chains.

\begin{center}
\small
\renewcommand{\arraystretch}{1.15}
\begin{tabular}{r@{\qquad}l@{\qquad}c}
\toprule
$N$ & $F$ & $a$\\
\midrule
$2598660833$ & $709\cdot2437$ & $3$\\
$P$ & $27953\cdot2598660833$ & $15$\\[1mm]
$114281$ & $2857$ & $6$\\
$10285291$ & $114281$ & $2$\\
$204710635813423$ & $6367\cdot10285291$ & $3$\\
$Q$ & $204710635813423$ & $3$\\
\bottomrule
\end{tabular}
\end{center}

Successive applications of Lemma~\ref{lem:pocklington} prove every
non-base entry prime, in particular $P$ and $Q$.
\end{proof}

\medskip
\noindent
\textbf{Exact verification.}
Every arithmetic assertion above is a finite computation in exact
integer arithmetic on the data printed in this paper, and can be redone
from the paper alone: multiply the seven factors
in~\eqref{eq:factorization} and compare the product with $2^{170}-1$;
confirm the five small factors of~\eqref{eq:factorization} and the five
base primes $709,2437,2857,6367,27953$ by trial division; for each row
of the table above check that $F\mid N-1$, that $F^2>N$, and that the
displayed witness $a$ satisfies the two conditions of
Lemma~\ref{lem:pocklington} at every prime $q\mid F$; compute
$\gcd(2^{170}-1,170)$ and the orders $\ord_5(2)=4$ and
$\ord_{17}(2)=8$; check that the squarefree products of three distinct
primes below $170$ are exactly the thirteen values displayed in the
proof of Theorem~\ref{thm:main}, that the three stated exclusions
remove all but $\{70,130,154,165\}$, and that the cyclotomic
factorizations of $\Phi_{35}(2)$, $\Phi_{130}(2)$ and $\Phi_{11}(2)$
displayed there are correct.

The computation was also re-run from scratch and independently, on a
separate machine, by a program that takes the numbers claimed here as
its input and recomputes everything else: all $36$ of its checks agreed
with the statements above. That re-run went beyond the list just given
in three ways: it proved $P$ and $Q$ prime a second time by Lucas
certificates built independently of the Pocklington chains, it
factored $2^n-1$ completely for each of the four exponents below $170$
that satisfy both hypotheses, and it confirmed by mutation that the
same predicates report failure when the exhibited factorization or
either large factor is perturbed. The program and its transcript
accompany this note. Two things the
re-run did not settle. A supplementary sweep of the bound of
Lemma~\ref{lem:cyclotomic} was carried out only for $n\le100$, because
factoring $2^n-1$ for prime exponents near $170$ was outside its
budget; that sweep is an extra, and the bound is proved above rather
than computed. And the bibliographic attributions in
Remark~\ref{rem:scope} are not checkable by machine; they were
confirmed by reading the sources. An independent re-run that agrees is
evidence that the computation is right; the proofs above, not the
program, are what is offered as the argument.

\begin{thebibliography}{9}

\bibitem{BLSTW}
J.~Brillhart, D.~H. Lehmer, J.~L. Selfridge, B.~Tuckerman, and
S.~S. Wagstaff,~Jr.,
\emph{Factorizations of $b^n\pm1$, $b=2,3,5,6,7,10,11,12$ up to high
powers}, 3rd ed., Contemp. Math. \textbf{22},
Amer. Math. Soc., Providence, RI, 2002.
\href{https://doi.org/10.1090/conm/022}{doi:10.1090/conm/022}.

\bibitem{DHS}
S.~Dolfi, R.~Hafezieh, and P.~Spiga,
\emph{On the structure of the character degree graphs having diameter three},
J. Algebra \textbf{685} (2026), 337--360.
\href{https://doi.org/10.1016/j.jalgebra.2025.08.002}
{doi:10.1016/j.jalgebra.2025.08.002};
preprint \href{https://arxiv.org/abs/2402.19335}{arXiv:2402.19335v1}
[math.GR], 2024. Conjecture~4.3 is quoted above from the preprint.

\bibitem{LLRS}
J.~Laubacher, M.~L. Lewis, L.~Ravaglia, and A.~Summers,
\emph{Constructing solvable groups whose character degree graphs
generalize the bowtie},
\href{https://arxiv.org/abs/2608.19374}{arXiv:2608.19374v1}
[math.GR], 2026.

\bibitem{OEIS}
OEIS Foundation Inc.,
\emph{The On-Line Encyclopedia of Integer Sequences},
Sequence A046800, ``Number of distinct prime factors of $2^n-1$,''
accessed August 22, 2026.
\href{https://oeis.org/A046800}{oeis.org/A046800}.
The terms displayed on the sequence's page stop at $n=101$; the value
at $n=170$ is in the linked b-file, which tabulates $a(n)$ for
$n=0,\dots,1206$.

\end{thebibliography}

\end{document}